% CALIBRATION DIAGNOSTICS (for Supplement §9.4, after Table S16)
% Generated: 2026-10-09

\subsection{Calibration Diagnostics}
\label{sec:calibration_diagnostics}

\paragraph{Structural validation.} The calibration scenario applies a proportional scaling factor $k=1.2131$ (derived from unrounded model outputs: $k = 0.491 / 0.404749$) to the pembrolizumab-arm $r_{LR}$ ratio curve while preserving DMFS and OS curves. Structural integrity checks over all 480 months confirm:
\begin{itemize}
    \item RFS $\leq$ DMFS: 0 violations (max excess: $0.00 \\times 10^{-10}$)
    \item DMFS $\leq$ OS: 0 violations (max excess: $0.00 \\times 10^{-10}$)
    \item RFS monotonicity: 0 non-monotonic points
    \item DMFS and OS unchanged: max $|$\Delta$|$ $< 10^{-15}$
\end{itemize}

\paragraph{Landmark fit comparison.} Table~S16 reports the calibration anchored to 49.1\% RFS at 60 months. The exact scale factor $k=1.2131$ yields RFS(60)=49.10\% (exactly anchored). Key timepoints:

\begin{table}[H]
\centering
\caption{Calibration Fit Quality at Key Timepoints}
\begin{tabular}{lrrrr}
\toprule
Month & Observed & Base RFS & Calibrated RFS & Residual (fitted $-$ observed) \\
\midrule
18 & 62.2\% & 66.45\% & 76.31\% & +4.25pp $\rightarrow$ +14.11pp \\
24 & 60.0\% & 61.39\% & 72.50\% & +1.39pp $\rightarrow$ +12.50pp \\
36 & 55.6\% & 52.92\% & 64.20\% & $-$2.68pp $\rightarrow$ +8.60pp \\
48 & 49.1\% & 46.09\% & 55.91\% & $-$3.01pp $\rightarrow$ +6.81pp \\
60 & 49.1\% & 40.47\% & 49.10\% & $-$8.63pp $\rightarrow$ 0.00pp \\
\bottomrule
\end{tabular}
\end{table}

Fit quality (observed timepoints): Base case MAE=3.99pp, RMSE=4.70pp; Calibrated MAE=8.40pp, RMSE=9.75pp. The calibration improves 60-month fit at the cost of deterioration at earlier timepoints.

\paragraph{Clipping diagnostics.} The scaling would produce $r_{LR}>1.0$ in 31 early months (months 0--30); these are clipped to $r_{LR}=1.0$. This affects LR state occupancy in the pembrolizumab arm: at month 60, LR occupancy decreases from 13.4\% (base) to 4.8\% (calibrated), reflecting improved RFS.

\paragraph{Economic validation.} Results from the same model run (not mixed sources):
\begin{itemize}
    \item $\Delta$LY: 4.82 (unchanged, OS preserved)
    \item $\Delta$QALY: 3.86 (vs 4.05 base)
    \item $\Delta$Cost: \$197,299 (vs \$161,947 base)
    \item ICER: \$51,053/QALY (reported in Table~S16)
    \item ICER increase: +27.7\% vs base case
\end{itemize}

\paragraph{Interpretation.} This calibration is a \textbf{local anchor stress test}, not a globally better-fitting curve. It improves the 60-month RFS fit (largest model residual: $-8.6$pp) while maintaining structural validity, but trades off fit quality at 18--48 months. The ICER increases by 27.6\%, demonstrating that the base-case conclusion is robust to this specific fitting discrepancy.
